\documentclass[12pt,a4paper]{article}
\usepackage[UTF8]{ctex}
\usepackage{geometry}
\usepackage{listings}
\usepackage{xcolor}
\usepackage{hyperref}
\usepackage{amsmath}
\usepackage{algorithm}
\usepackage{algpseudocode}

\geometry{margin=2.5cm}

\lstset{
    language=Python,
    basicstyle=\ttfamily\small,
    keywordstyle=\color{blue}\bfseries,
    commentstyle=\color{green!60!black},
    stringstyle=\color{red},
    numbers=left,
    numberstyle=\tiny\color{gray},
    stepnumber=1,
    numbersep=5pt,
    frame=single,
    breaklines=true,
    showstringspaces=false,
    tabsize=4
}

\title{Halo 系统代码文档：Parser 模块}
\author{代码分析文档}
\date{\today}

\begin{document}

\maketitle

\tableofcontents
\newpage

\section{概述}

\texttt{parser.py} 模块负责从 YAML 配置文件解析并构建智能体工作流的有向无环图（DAG）。它是整个系统的入口点，将声明式的工作流描述转换为可执行的 \texttt{Operator} 对象图。

\section{模块功能}

\begin{itemize}
    \item 加载和解析 YAML 配置文件
    \item 验证配置的完整性和正确性
    \item 构建 \texttt{Operator} 对象图（DAG）
    \item 计算每个 OP 到终点的最长距离（用于调度优化）
    \item 检测图中的循环依赖
\end{itemize}

\section{YAML 配置格式}

\subsection{基本结构}

\begin{lstlisting}
start_ops:
  - op0          # 起始节点列表
end_ops:
  - op2          # 终止节点列表

ops:
  op0:
    model: meta-llama/Llama-3.1-8B-Instruct
    prompt: "Please answer the question."
    max_tokens: 512
    input_ops: []           # 空列表表示起始节点
    output_ops:
      - op1
    
  op1:
    model: meta-llama/Llama-3.1-8B-Instruct
    prompt: "Please refine the answer."
    max_tokens: 1024
    input_ops:
      - op0      # 依赖 op0 的输出
    output_ops:
      - op2
    
  op2:
    model: meta-llama/Llama-3.1-8B-Instruct
    prompt: "Please provide final answer."
    max_tokens: 512
    input_ops:
      - op1
    output_ops: []          # 空列表表示终止节点
\end{lstlisting}

\subsection{字段说明}

\begin{description}
    \item[\texttt{start_ops}] 必需，非空列表，指定工作流的入口节点
    \item[\texttt{end_ops}] 必需，非空列表，指定工作流的出口节点
    \item[\texttt{ops}] 必需，字典，键为 OP ID，值为 OP 配置
    \item[\texttt{model}] 必需，模型标识符
    \item[\texttt{prompt}] 可选，系统提示词
    \item[\texttt{input_ops}] 可选，前驱节点列表，默认为空列表
    \item[\texttt{output_ops}] 可选，后继节点列表，默认为空列表
    \item[\texttt{max_tokens}] 可选，最大生成 token 数，默认 256
    \item[\texttt{temperature}] 可选，采样温度，默认 0.7
    \item[\texttt{top_p}] 可选，Top-p 采样参数，默认 0.9
    \item[\texttt{keep_cache}] 可选，是否保留 KV cache，默认自动推断
\end{description}

\section{核心函数}

\subsection{load\_config}

\begin{lstlisting}
def load_config(config_path: str) -> dict:
    """Load YAML config as dict."""
    with open(config_path, "r") as f:
        return yaml.safe_load(f)
\end{lstlisting}

\begin{description}
    \item[功能] 从文件路径加载 YAML 配置文件并解析为 Python 字典
    \item[参数] \texttt{config_path}：YAML 文件的路径
    \item[返回] 解析后的配置字典
    \item[实现] 使用 \texttt{yaml.safe\_load} 安全地解析 YAML，避免执行任意代码
\end{description}

\subsection{build\_ops\_from\_config}

\subsubsection{函数签名}

\begin{lstlisting}
def build_ops_from_config(config: dict) -> Tuple[
    Dict[str, Operator],  # ops: op_id -> Operator
    List[Operator],       # start_ops
    List[Operator],      # end_ops
    Set[str]             # models: 所有使用的模型名称
]:
\end{lstlisting}

\subsubsection{功能概述}

这是 parser 模块的核心函数，负责将配置字典转换为可执行的 \texttt{Operator} 对象图。

\subsubsection{执行流程}

\begin{enumerate}
    \item \textbf{验证必需字段}
    \begin{lstlisting}
    conf_ops = config.get("ops")
    if not isinstance(conf_ops, dict) or not conf_ops:
        raise ValueError("Config must contain a non-empty 'ops' mapping.")
    start_keys = config.get("start_ops")
    end_keys = config.get("end_ops")
    if not isinstance(start_keys, list) or not start_keys:
        raise ValueError("'start_ops' must be a non-empty list of op ids.")
    if not isinstance(end_keys, list) or not end_keys:
        raise ValueError("'end_ops' must be a non-empty list of op ids.")
    \end{lstlisting}
    
    \item \textbf{创建 Operator 对象}
    \begin{lstlisting}
    ops: Dict[str, Operator] = {oid: Operator(id=oid) for oid in conf_ops.keys()}
    \end{lstlisting}
    为每个 OP ID 创建一个空的 \texttt{Operator} 对象。
    
    \item \textbf{验证引用和必需字段}
    \begin{lstlisting}
    for oid, spec in conf_ops.items():
        if "model" not in spec:
            raise ValueError(f"Op '{oid}' is missing required field 'model'.")
        in_ids = spec.get("input_ops", []) or []
        out_ids = spec.get("output_ops", []) or []
        for rid in in_ids + out_ids:
            if rid not in ops:
                raise ValueError(f"Op '{oid}' references unknown op '{rid}' in inputs/outputs.")
    \end{lstlisting}
    检查：
    \begin{itemize}
        \item 每个 OP 必须有 \texttt{model} 字段
        \item \texttt{input_ops} 和 \texttt{output_ops} 中引用的 OP ID 必须存在
    \end{itemize}
    
    \item \textbf{建立依赖关系并配置模型}
    \begin{lstlisting}
    for oid, op in ops.items():
        spec = conf_ops[oid]
        input_ids = spec.get("input_ops", []) or []
        output_ids = spec.get("output_ops", []) or []
        
        # 建立双向链接
        op.input_ops = [ops[k] for k in input_ids]
        op.output_ops = [ops[k] for k in output_ids]
        
        model = spec["model"]
        models.add(model)
        
        # keep_cache 推断逻辑
        explicit_keep = spec.get("keep_cache", None)
        inferred_keep = any(conf_ops[oid2]["model"] == model for oid2 in output_ids)
        op.keep_cache = bool(explicit_keep) if explicit_keep is not None else inferred_keep
    \end{lstlisting}
    
    \textbf{关键逻辑：keep\_cache 推断}
    \begin{itemize}
        \item 如果配置中显式指定了 \texttt{keep\_cache}，则使用该值
        \item 否则，如果该 OP 的任意后继节点使用相同模型，则自动推断为 \texttt{True}
        \item 这有助于优化：如果下游 OP 使用相同模型，保留 cache 可以避免重复计算前缀
    \end{itemize}
    
    \item \textbf{创建 ModelConfig 对象}
    \begin{lstlisting}
    op.model_config = ModelConfig(
        model_name=model,
        system_prompt=spec.get("prompt"),
        temperature=spec.get("temperature", 0.7),
        top_p=spec.get("top_p", 0.9),
        max_tokens=spec.get("max_tokens", 256),
        max_batch_size=spec.get("max_batch_size", torch.inf),
        dtype=spec.get("dtype", "bfloat16"),
        quantization=spec.get("quantization", None),
        lora_config=spec.get("lora_config", None),
        max_model_len=spec.get("max_model_len", None),
        min_tokens=spec.get("min_tokens", 0),
        use_chat_template=spec.get("use_chat_template", True),
    )
    \end{lstlisting}
    
    \item \textbf{初始化运行时字段}
    \begin{lstlisting}
    op.data_parallel = False
    op.is_duplicate = False
    op.duplicate_info = None
    op.main_op = None
    op.max_distance = -1  # 稍后由 _compute_max_distances 计算
    \end{lstlisting}
    
    \item \textbf{解析 start\_ops 和 end\_ops}
    \begin{lstlisting}
    for k in start_keys + end_keys:
        if k not in ops:
            raise ValueError(f"Unknown op id '{k}' referenced in start_ops/end_ops.")
    start_ops = [ops[k] for k in start_keys]
    end_ops = [ops[k] for k in end_keys]
    \end{lstlisting}
    
    \item \textbf{计算最长距离}
    \begin{lstlisting}
    _compute_max_distances(ops, end_ops)
    \end{lstlisting}
    为每个 OP 计算到任意终点 OP 的最长路径长度，用于调度优化。
    
    \item \textbf{返回结果}
    \begin{lstlisting}
    return ops, start_ops, end_ops, models
    \end{lstlisting}
\end{enumerate}

\subsection{build\_from\_path}

\begin{lstlisting}
def build_from_path(config_path: str) -> Tuple[
    Dict[str, Operator],
    List[Operator],
    List[Operator],
    Set[str]
]:
    """Convenience wrapper: load + build in one call."""
    return build_ops_from_config(load_config(config_path))
\end{lstlisting}

\begin{description}
    \item[功能] 便捷函数，一次性完成加载和构建
    \item[参数] \texttt{config_path}：YAML 文件路径
    \item[返回] 与 \texttt{build\_ops\_from\_config} 相同
    \item[用途] 简化调用，避免手动调用两个函数
\end{description}

\section{最长距离计算}

\subsection{\_compute\_max\_distances}

\subsubsection{功能}

为每个 OP 计算到任意终点 OP 的最长路径长度（边数），同时检测循环依赖。

\subsubsection{算法：DFS + 记忆化}

\begin{algorithm}
\caption{计算最长距离到终点}
\begin{algorithmic}[1]
\Procedure{DFS}{op, end\_set, memo, visiting}
    \If{op in memo}
        \State \Return memo[op]
    \EndIf
    \If{op in visiting}
        \State \textbf{raise} "Cycle detected"
    \EndIf
    \State visiting.add(op)
    \If{op in end\_set}
        \State memo[op] = 0
        \State visiting.remove(op)
        \State \Return 0
    \EndIf
    \State best = -1
    \For{child in op.output\_ops}
        \State d = DFS(child, end\_set, memo, visiting)
        \If{d != -1}
            \State best = max(best, d + 1)
        \EndIf
    \EndFor
    \State memo[op] = best
    \State visiting.remove(op)
    \State \Return best
\EndProcedure
\end{algorithmic}
\end{algorithm}

\subsubsection{实现细节}

\begin{lstlisting}
def _compute_max_distances(ops: Dict[str, Operator], end_ops: List[Operator]) -> None:
    """
    For each op, compute the longest distance (in edges) to ANY end-op.
    If an op cannot reach any end-op, its distance is -1.
    Detect cycles and raise ValueError if found.
    """
    end_set = set(end_ops)
    memo: Dict[Operator, int] = {}
    visiting: Set[Operator] = set()  # 用于循环检测
    
    def dfs(op: Operator) -> int:
        if op in memo:
            return memo[op]
        if op in visiting:
            raise ValueError("Cycle detected in the op graph.")
        visiting.add(op)
        
        if op in end_set:
            memo[op] = 0
            visiting.remove(op)
            return 0
        
        best = -1
        for child in op.output_ops:
            d = dfs(child)
            if d != -1:  # 可达终点
                best = max(best, d + 1)
        
        memo[op] = best
        visiting.remove(op)
        return best
    
    # 为所有 OP 计算距离
    for op in ops.values():
        op.max_distance = dfs(op)
\end{lstlisting}

\subsubsection{关键点}

\begin{itemize}
    \item \textbf{记忆化}：避免重复计算，时间复杂度 $O(V + E)$
    \item \textbf{循环检测}：使用 \texttt{visiting} 集合检测后向边，如果访问到正在访问的节点，说明存在环
    \item \textbf{不可达处理}：如果 OP 无法到达任何终点，距离为 -1
    \item \textbf{最长路径}：使用 \texttt{max} 选择所有后继节点中的最长路径
\end{itemize}

\subsubsection{用途}

\texttt{max\_distance} 字段在调度器中被广泛使用：
\begin{itemize}
    \item \textbf{优先级排序}：距离终点越远的 OP 通常应该优先执行（尽早开始长链）
    \item \textbf{调度决策}：帮助调度器决定哪些 OP 应该先分配到设备
    \item \textbf{负载均衡}：考虑剩余工作量，更好地平衡各设备的负载
\end{itemize}

\section{错误处理}

\subsection{可能的错误}

\begin{enumerate}
    \item \textbf{缺少必需字段}
    \begin{itemize}
        \item \texttt{ops} 为空或不存在
        \item \texttt{start\_ops} 或 \texttt{end\_ops} 为空
        \item 某个 OP 缺少 \texttt{model} 字段
    \end{itemize}
    
    \item \textbf{引用错误}
    \begin{itemize}
        \item \texttt{input\_ops} 或 \texttt{output\_ops} 中引用了不存在的 OP ID
        \item \texttt{start\_ops} 或 \texttt{end\_ops} 中引用了不存在的 OP ID
    \end{itemize}
    
    \item \textbf{循环依赖}
    \begin{itemize}
        \item 图中存在有向环（通过 DFS 检测）
        \item 抛出 \texttt{ValueError("Cycle detected in the op graph.")}
    \end{itemize}
\end{enumerate}

\section{使用示例}

\subsection{基本用法}

\begin{lstlisting}
from halo.parser import build_from_path

# 从 YAML 文件构建工作流图
ops, start_ops, end_ops, models = build_from_path("templates/adv_reason_3.yaml")

# 访问特定 OP
op0 = ops["op0"]
print(f"OP0 使用模型: {op0.model_config.model_name}")
print(f"OP0 的前驱: {[op.id for op in op0.input_ops]}")
print(f"OP0 的后继: {[op.id for op in op0.output_ops]}")
print(f"OP0 到终点的最长距离: {op0.max_distance}")
\end{lstlisting}

\subsection{验证图结构}

\begin{lstlisting}
# 检查所有 OP 是否可达终点
for op_id, op in ops.items():
    if op.max_distance == -1:
        print(f"警告: {op_id} 无法到达任何终点节点")
    else:
        print(f"{op_id} 到终点的最长距离: {op.max_distance}")
\end{lstlisting}

\section{总结}

\texttt{parser.py} 模块是 Halo 系统的配置解析层，主要功能包括：

\begin{itemize}
    \item \textbf{配置解析}：从 YAML 文件加载工作流定义
    \item \textbf{图构建}：将配置转换为 \texttt{Operator} 对象的有向无环图
    \item \textbf{依赖建立}：建立 OP 之间的输入输出依赖关系
    \item \textbf{自动推断}：自动推断 \texttt{keep\_cache} 等优化参数
    \item \textbf{距离计算}：计算每个 OP 到终点的最长距离，用于调度优化
    \item \textbf{错误检测}：检测配置错误和循环依赖
\end{itemize}

该模块的输出（\texttt{ops}、\texttt{start\_ops}、\texttt{end\_ops}）是后续调度器和优化器的输入，是整个系统的基础。

\end{document}
